\chapter*{คำนำ}
\addcontentsline{toc}{chapter}{คำนำ}

การบริหารจัดการสวนทุเรียนยุคดิจิทัลจำเป็นต้องอาศัยการตรวจติดตามสุขภาพทางสรีรวิทยาและโรคพืชอย่างแม่นยำและทันท่วงที ในอดีตการตรวจหาอาการโรคใบพืชและการประเมินความสมบูรณ์ของธาตุอาหารจำเป็นต้องพึ่งพาความเชี่ยวชาญเฉพาะบุคคลของเกษตรกรหรือการส่งตัวอย่างใบไม้เข้าตรวจวิเคราะห์ในห้องปฏิบัติการทางเคมี ซึ่งใช้ระยะเวลานานและมีค่าใช้จ่ายสูง การเกิดโรคระบาดรุนแรง เช่น โรคใบไหม้ไฟทอปธอร่า (\textit{Phytophthora palmivora}) หรือโรคราใบติด (\textit{Rhizoctonia solani}) จึงมักสร้างความเสียหายอย่างมีนัยสำคัญก่อนที่การรักษาจะเริ่มต้นขึ้น

หน่วยวิเคราะห์-ติดตามคุณภาพดินและสภาพอากาศเชิงพื้นที่ คณะวิทยาศาสตร์และเทคโนโลยี มหาวิทยาลัยราชภัฏรำไพพรรณี จึงได้บูรณาการองค์ความรู้ด้านฟิสิกส์เชิงแสงและสเปกโทรสโกปีเสมือน ร่วมกับโมเดลการเรียนรู้เชิงลึกแบบ Dual-Head Convolutional Neural Network ออกแบบและพัฒนาแอปพลิเคชันมือถือ \textbf{DurianLeaf AI (PlantHealth)} ที่สามารถตรวจวินิจฉัยโรคพืชควบคู่กับการประเมินภาวะโภชนาการพืช (ธาตุอาหารหลัก NPK ธาตุอาหารรอง และจุลธาตุ) ผ่านกล้องถ่ายภาพสดแบบเรียลไทม์

จุดเด่นสำคัญของระบบคือ แถบเมนูปรับขนาดและรูปทรงของพื้นที่ตรวจจับ (Interactive ROI Floating Toolbar) ที่มีแถบแสดงเฉดสีตามเกณฑ์มาตรฐานทางวิชาการเรืองแสงเชื่อมโยงกับค่าการวิเคราะห์สด พร้อมทั้งกลไกการเรียนรู้เพิ่มเติมระดับเครื่อง (On-Device Continual Few-Shot Learning) ที่เปิดโอกาสให้ผู้ใช้งานสามารถสอนอาการโรคใหม่หรือสายพันธุ์ทุเรียนใหม่เข้าสู่หน่วยความจำของโมเดลได้ทันทีโดยไม่ต้องคอมไพล์โปรแกรมใหม่

คณะผู้จัดทำหวังเป็นอย่างยิ่งว่าคู่มือปฏิบัติการฉบับนี้ จะเป็นประโยชน์แก่นักวิจัย เกษตรกรผู้ปลูกทุเรียน ตลอดจนคณาจารย์และนักศึกษา เพื่อขับเคลื่อนเทคโนโลยีการเกษตรแม่นยำของประเทศไทยสู่ระดับสากลต่อไป

\vspace{1.5cm}

\begin{flushright}
  \textbf{ผู้ช่วยศาสตราจารย์ ดร.ชีวะ ทัศนา}\\
  หน่วยวิเคราะห์-ติดตามคุณภาพดินและสภาพอากาศเชิงพื้นที่\\
  คณะวิทยาศาสตร์และเทคโนโลยี มหาวิทยาลัยราชภัฏรำไพพรรณี
\end{flushright}
